\section{Reduction to approximation and rounding}\label{sec:reduction}
For a vector $X=(X_0,\ldots,X_{N-1})\in[-1,1]^N$, write
\[
 Q_X(t)=\sum_{k=0}^{N-1}X_k\e(kt),\qquad
 \mu(X)=\frac12\sum_{k=0}^{N-1}(1-|X_k|).
\]
The first proposition constructs vectors with small Fourier maximum and
nearly maximal energy. Its proof occupies Sections~\ref{sec:packing}--\ref{sec:sampling}.

\begin{proposition}[Almost-sign approximation]\label{prop:almost-signs}
For every $0<\delta<1/20$, there are vectors
$X_N=(X_{N,0},\ldots,X_{N,N-1})\in[-1,1]^N$, one for each integer
$N\ge1$, such that
\begin{align}
 \limsup_{N\to\infty}\frac{\norm{Q_{X_N}}_\infty}{\sqrt N}
 &\le K_\delta,
 &K_\delta&=\sqrt{\frac{(1+\delta)^3}{1-\delta}},
 \label{eq:almost-maximum}\\
 \liminf_{N\to\infty}\frac1N\sum_{k=0}^{N-1}X_{N,k}^2
 &\ge1-7\delta.\label{eq:almost-energy}
\end{align}
\end{proposition}

The next lemma converts an arbitrary relaxed vector to signs. Its
constant does not depend on the vector or its length. We prove it in
Section~\ref{sec:rounding}.

\begin{lemma}[Rounding with a small defect]\label{lem:rounding}
There is an absolute constant $C$ such that for every $N\ge1$ and every
$X\in[-1,1]^N$, some signs $\varepsilon_0,\ldots,\varepsilon_{N-1}$
satisfy
\begin{equation}\label{eq:defect-rounding}
 \max_{t\in\T}\left|\sum_{k=0}^{N-1}(\varepsilon_k-X_k)\e(kt)\right|
 \le C\left(1+\sqrt{\mu(X)\log\frac{80N}{\mu(X)}}\right).
\end{equation}
The square-root term is defined to be zero when $\mu(X)=0$; in that case
$X$ already has sign entries and the left side may be zero.
\end{lemma}

\begin{proof}[Proof of Theorem~\ref{thm:main}]
Fix $0<\delta<1/20$ and choose the family from
Proposition~\ref{prop:almost-signs}. For all sufficiently large $N$,
\eqref{eq:almost-energy} gives
\[
 \frac1N\sum_kX_{N,k}^2\ge1-8\delta.
\]
Since $1-|x|\le1-x^2$ for $|x|\le1$, it follows that
$\mu(X_N)\le4\delta N$. The function
$u\mapsto u\log(80N/u)$ is increasing on $0<u\le N/2$.
Applying Lemma~\ref{lem:rounding} and then
\eqref{eq:almost-maximum}, we obtain
\[
 \limsup_{N\to\infty}m_N
 \le K_\delta+2C\sqrt{\delta\log(20/\delta)}.
\]
First the limit in $N$ was taken with every auxiliary choice depending
on $\delta$ fixed. Now let $\delta\downarrow0$. The right side tends to
one, and the lower bound $m_N\ge1$ follows from Parseval. This proves
the limit through all positive integer lengths.
\end{proof}
